% =========================================================
% Families of Sets as Ordered Sets
% =========================================================

\section{Families of Sets}
\label{sec:families-of-sets}

\begin{remark*}[Toolkit: Families of Sets]
The power set of any set is naturally ordered by inclusion.  Every
collection of subsets --- whether arising from algebra, topology, or logic
--- inherits this order.  The key observation is that subsets correspond
exactly to predicates via their characteristic functions, and that this
correspondence is an order-isomorphism.
\end{remark*}

\begin{remark*}[Families as ordered sets]
Let $X$ be a set.  The \emph{power set} $\mathcal{P}(X)$, consisting of
all subsets of $X$, is ordered by inclusion: for $A, B \in \mathcal{P}(X)$,
\[
A \leq B \;\Longleftrightarrow\; A \subseteq B.
\]
Any subcollection of $\mathcal{P}(X)$ inherits this order.  Such a
collection is called a \emph{family of sets}.

Families of sets arise naturally when $X$ carries additional structure.
If $X$ is a group $G$, the set $\operatorname{Sub}(G)$ of all subgroups
and the set $\operatorname{NSub}(G)$ of all normal subgroups are each
ordered by inclusion.  If $X$ is a vector space $V$, the set
$\operatorname{Sub}(V)$ of all subspaces is ordered by inclusion.  If $X$
is a ring $R$, the sets of all subrings and of all ideals of $R$ are
similarly ordered.  If $(X, \mathcal{T})$ is a topological space, the
collections of open sets, closed sets, and clopen sets are each families
of sets ordered by inclusion.

A particularly important example is the family $\mathcal{O}(P)$ of all
\hyperref[def:down-set]{down-sets} of a partially ordered set $P$,
ordered by inclusion.
\end{remark*}

\begin{remark*}[Predicate formulation]
The ordered set $(\mathcal{P}(X), \subseteq)$ admits an equivalent
formulation in terms of predicates.  A \emph{predicate} on $X$ is a
function
\[
p : X \to \{T, F\},
\]
where $T$ and $F$ denote truth values.  Let $\mathcal{P}^*(X)$ denote
the set of all predicates on $X$, ordered by implication:
\[
p \leq q
\;\Longleftrightarrow\;
\{x \in X : p(x) = T\} \subseteq \{x \in X : q(x) = T\}.
\]

Define the map
\[
\Phi : \mathcal{P}^*(X) \to \mathcal{P}(X),
\qquad
\Phi(p) \coloneqq \{x \in X : p(x) = T\}.
\]
\end{remark*}

\begin{propositionbox}{Proposition}
\begin{proposition}[Predicates and Power Sets Are Order-Isomorphic]
\label{prop:predicates-power-sets-iso}
The map $\Phi : (\mathcal{P}^*(X), \leq) \to (\mathcal{P}(X), \subseteq)$
is an \hyperref[def:order-isomorphism]{order-isomorphism}.
\hyperref[prf:predicates-power-sets-iso]{Go to proof.}
\end{proposition}
\end{propositionbox}

\begin{remark*}[Standard quantified statement]
\[
\operatorname{OrderIso}(\Phi, \mathcal{P}^*(X), \mathcal{P}(X)).
\]
\end{remark*}

\begin{remark*}[Interpretation]
Every subset of $X$ can be represented as the truth set of a predicate,
and every predicate determines a subset.  The correspondence is bijective
and order-preserving in both directions: $p \leq q$ if and only if the
truth set of $p$ is contained in the truth set of $q$, which is exactly
the inclusion order on $\mathcal{P}(X)$.  This identification is
fundamental in logic, where sets and their characteristic predicates are
treated as interchangeable, and in computer science, where sets are
often implemented via their membership tests.
\end{remark*}

\begin{dependencies}
\hyperref[def:order-isomorphism]{Order-isomorphism},
\hyperref[def:lattice-order-partial-order]{partial order}.
\end{dependencies}
